\cgeentry{cge_structure}{Structure of the system: structural singularity and blocks}

\cgesec{Description}

\cmd{cge_structure} (from version 0.35.0) analyses the structure of the system
that \manref{cge_solve} solves under the closure: which free variables each
equation contains, from the nonzeros of its exact Jacobian at a point. A
maximum matching assigns variables to equations. If it covers every equation
and every variable, the system is structurally nonsingular; if not, it is
singular whatever its values, and the Dulmage--Mendelsohn decomposition names
its parts: the over-determined part, equations that together contain fewer
free variables than they number, and the under-determined part, variables
that the equations cannot all determine. The square part is split into
blocks that can be solved in turn, each given the ones before it, as CONOPT
finds a model's pre-triangular, simultaneous and post-triangular parts.

\cgesec{Quick start}

The structure of the system under the current closure
\begin{stcode}
cge_structure
\end{stcode}
The full system at the base, with the twenty largest blocks
\begin{stcode}
cge_structure, base full top(20)
\end{stcode}

\cgesec{Syntax}

\begin{stsyntax}
\cmdb{cge_structure} [\texttt{,} \optarg{}{sim}{name} \opt{}{base} \opt{}{full} \optarg{}{top}{\#}]
\end{stsyntax}

\begin{opttable}
\optgroup{Main}
\optarg{}{sim}{name} & the system at simulation \argu{name}\\
\opt{}{base} & the system at the base\\
\opt{}{full} & the full system, with the definitions that the reduced one leaves out\\
\optarg{}{top}{\#} & the labels and blocks listed; default \texttt{10}\\
\end{opttable}

\cgesec{Options}

\begin{optdesc}
\item[\optarg{}{sim}{name}, \opt{}{base}] choose the point: a stored
simulation or the base. By default, the current point, with the current
shocks and closure, from which the next solve starts.

\item[\opt{}{full}] analyses the full system. By default, the system the
solver uses: the reduced one for a model from a SAM (section~\ref{sec:reduced}).

\item[\optarg{}{top}{\#}] limits the labels of each part and the blocks
listed.
\end{optdesc}

\cgesec{Remarks and examples}

A change of closure, a swap or a user equation can leave a system that is
square, as many free variables as equations, yet singular whatever its
values: some equations contain too few free variables between them, and some
variables appear in too few equations. Newton's method then meets a singular
Jacobian, or converges to a point that the system does not determine.
\cmd{cge_structure} finds such a system before it is solved and names the
equations and variables at fault.

The SAM is mapped and calibrated:
\stlog{struct_setup}

\cgesubsec{Example 1: the structure under a closure}

The reduced system has a maximum matching that covers it: it is structurally
nonsingular. Its equations fall into blocks that can be solved in turn; most
are single equations, and one block holds the simultaneous core of prices
and quantities:
\stlog{struct_default}

\cgesubsec{Example 2: a structurally singular system}

A user variable that no equation contains, and a user equation in the
numeraire alone, keep the system square but make it singular. The equation
is over-determined, with no free variable to determine; the variable is
under-determined, with no equation to determine it:
\stlog{struct_singular}

\cgesubsec{Example 3: a standalone model}

A standalone model (\manref{cge_model}) decomposes in the same way. Here
the first three equations can be solved one after the other, the next two
together, and the last given those:
\stlog{struct_standalone}

\cgesec{Stored results}

\cmd{cge_structure} stores the following in \cmd{r()}:
\begin{storedres}
\resgroup{Scalars}
\res{r(N_equations)} & equations of the system\\
\res{r(N_variables)} & free variables\\
\res{r(rank)} & structural rank, the size of a maximum matching\\
\res{r(singular)} & 1 if the system is structurally singular\\
\res{r(N_over_eq)} & equations of the over-determined part\\
\res{r(N_over_var)} & its variables\\
\res{r(N_under_eq)} & equations of the under-determined part\\
\res{r(N_under_var)} & its variables\\
\res{r(N_square)} & equations of the square part\\
\res{r(N_blocks)} & blocks of the square part\\
\res{r(largest_block)} & equations of the largest block\\
\res{r(N_singletons)} & blocks of one equation\\
\resgroup{Macros}
\res{r(over_equations)} & equations of the over-determined part\\
\res{r(over_variables)} & its variables\\
\res{r(under_equations)} & equations of the under-determined part\\
\res{r(under_variables)} & its variables\\
\end{storedres}

\cgesec{Methods and formulas}

The pattern is the set of nonzero entries of the exact Jacobian at the point
(chapter~\ref{chap:solution}), rows of the system and its free columns. A
maximum matching is found by assigning each row, in order of its number of
entries, to its first unmatched column, and then from each row left
unmatched by a depth-first search for an augmenting path. With $R_0$ the
unmatched rows and $C_0$ the unmatched columns, the over-determined part is
every row and column reached from $R_0$ by alternating paths (a row to its
columns, a column to the row matched to it), and the under-determined part
every column and row reached from $C_0$ (a column to its rows, a row to the
column matched to it) \citep{dulmage1958}. In the square part the rows form a
directed graph, row $i$ pointing to row $k$ when $i$ contains the column
matched to $k$; its strongly connected components, found by Tarjan's
algorithm \citep{tarjan1972}, are the blocks, in an order in which each
depends only on itself and on those before it.

\cgesec{Also see}

\begin{alsosee}
\seeentry{cge_listing}{Equation listing, infeasibility report and solution listing}
\seeentry{cge_closure}{Choose the model closure}
\seeentry{cge_solve}{Solve the model}
\seeentry{cge_model}{A standalone model, written whole in the user language}
\end{alsosee}
